\documentclass{beamer}

\mode<presentation> {
\usetheme{Singapore}
\setbeamercovered{transparent}

%\setbeamertemplate{footline} % To remove the footer line in all slides uncomment this line
%\setbeamertemplate{footline}[page number] % To replace the footer line in all slides with a simple slide count uncomment this line
%\setbeamertemplate{navigation symbols}{} % To remove the navigation symbols from the bottom of all slides uncomment this line
}

% \usepackage{ctex}
\usepackage{fontawesome}
\usepackage{graphicx} 
\usepackage{booktabs} 
\usepackage{xcolor}
\usepackage{calc}
\usepackage{amsmath}
\usepackage{amssymb}

\setbeamertemplate{caption}[numbered]

% \makeatletter
% \setbeamertemplate{frametitle}{
%     \ifbeamercolorempty[bg]{frametitle}{}{\nointerlineskip}%
%     \@tempdima=\textwidth%
%     \advance\@tempdima by\beamer@leftmargin%
%     \advance\@tempdima by\beamer@rightmargin%
%     \vspace*{1cm} %%%%%%%%%%%%% For example insert shift to right
%     \begin{beamercolorbox}[sep=0.3cm,center,wd=\the\@tempdima]{frametitle}
%         \usebeamerfont{frametitle}%
%         \vbox{}\vskip-1ex%
%         \if@tempswa\else\csname beamer@ftecenter\endcsname\fi%
%         {%
%             \ifx\insertframesubtitle\@empty%
%             \else%
%             {\usebeamerfont{framesubtitle}\usebeamercolor[fg]{framesubtitle}\insertframesubtitle\strut\par}%
%             \fi
%         }%
%         \vskip-1ex%
%         \if@tempswa\else\vskip-.3cm\fi% set inside beamercolorbox... evil here...
%     \end{beamercolorbox}%
% }
% \makeatother

\newcommand{\vect}[1]{\boldsymbol{#1}}
\newcommand{\matr}[1]{\mathbf{#1}}
\newcommand{\lnn}[1]{
  \ln\left(#1\right)
}

\newlength\dlf
\newcommand\alignedbox[2]{
  % #1 = before alignment
  % #2 = after alignment
  &
  \begingroup
  \settowidth\dlf{$\displaystyle #1$}
  \addtolength\dlf{\fboxsep+\fboxrule}
  \hspace{-\dlf}
  \fcolorbox{red}{yellow}{$\displaystyle #1 #2$}
  \endgroup
}

\pgfdeclareimage[height=0.7cm]{company-logo}{img/mediatek.png}
\logo{\pgfuseimage{company-logo}}

\author{\Huge{Tzu-Chun Hsu}\\[10pt]} 
\institute[Realtek Semiconductor Corp.] 
{
    \scriptsize{
        \href{tel:+886-987605719}{ \raisebox{-0.1\height}\faPhone\ \underline{+886-987605719} ~} 
        \href{mailto:vm3y3rmp40719@gmail.com}{\raisebox{-0.2\height}\faEnvelope\  \underline{vm3y3rmp40719@gmail.com}} \\~\\
        \href{https://www.linkedin.com/in/tzu-chun-hsu-ab4b3b188/}{\raisebox{-0.2\height}\faLinkedinSquare\ \underline{tzu-chun-hsu-ab4b3b188} ~}
        \href{https://github.com/Oscarshu0719}{\raisebox{-0.2\height}\faGithub\ \underline{Oscarshu0719}}
    }
}
\date{March 24, 2026} 

\begin{document}

\begin{frame}
\titlepage % Print the title page as the first slide
\end{frame}

% \begin{frame}
% \frametitle{Table of contents} 
% \tableofcontents 
% \end{frame}

%----------------------------------------------------------------------------------------
%	WORK EXPERIENCE SLIDES
%----------------------------------------------------------------------------------------

\section{Work Experience}
\begin{frame}
    \frametitle{Work Experience}
    \begin{block}{Camera ISP Firmware Engineer}
        \begin{overlayarea}{\textwidth}{0.7\textheight}
            Realtek Semiconductor Corp.\\
            Hsinchu, Taiwan \hfill 02 2025 – present
            \begin{enumerate}
                \item \textbf{ISP Development \& Algorithm Optimization.}
                \item \textbf{Software Architecture \& Development.}
                \item \textbf{Test Automation \& Verification.}
            \end{enumerate}
        \end{overlayarea}
    \end{block}
\end{frame}

\begin{frame}
    \frametitle{Work Experience}
    \begin{block}{Camera ISP Firmware Engineer}
        \begin{overlayarea}{\textwidth}{0.7\textheight}
            Realtek Semiconductor Corp.\\
            Hsinchu, Taiwan \hfill 02 2025 – present
            \begin{enumerate}
                \item<1> \textbf{ISP Development \& Algorithm Optimization:}
                \only<1>{\begin{itemize}
                    \item Developed high-resolution image stitching solution for smart glasses.
                    \item Aligned ISP pipeline for snapshot and streaming capture to ensure consistent output.
                    \item Migrated ISP engine for smart glasses from Linux to FreeRTOS.
                    \item Implemented ISP pipeline features for smart glasses with V4L2 and interrupt controls and optimized SDK flow.
                \end{itemize}}
                \item<2> \textbf{Software Architecture \& Development:}
                \only<2>{\begin{itemize}
                    \item Refactored core CModel architecture and consolidated sub-applications.
                    \item Redesigned more detailed and readable image format table.
                    \item Implemented thread and cross-process synchronization for Windows.
                \end{itemize}}
                \item<3> \textbf{Test Automation \& Verification:}
                \only<3>{\begin{itemize}
                    \item Designed and optimized automated verification systems for ISP/IPU (FPGA/ASIC) hardware platforms.
                    \item Migrated communication protocol from JRPC to protobuf for the Linux-to-FreeRTOS migration.
                \end{itemize}}
            \end{enumerate}
        \end{overlayarea}
    \end{block}
\end{frame}

%----------------------------------------------------------------------------------------
%	EDUCATION SLIDES
%----------------------------------------------------------------------------------------

\section{Education}
\begin{frame}
    \frametitle{Education}
    \begin{block}{Zhejiang University} % #47 QS 2025.
        Bachelor of Engineering in Computer Science and Technology\\
        Hangzhou, China \hfill 09 2016 – 07 2020
    \end{block}
    \begin{block}{National Yang Ming Chiao Tung University}
        Master of Science in Computer Science and Engineering\\
        Hsinchu, Taiwan \hfill 09 2022 – 02 2025
    \end{block}
\end{frame}

%----------------------------------------------------------------------------------------
%	PROJECTS SLIDES
%----------------------------------------------------------------------------------------

\section{Projects}
\begin{frame}[t]
    \frametitle{Projects}
    \begin{block}{\href{https://docs.google.com/presentation/d/1ge2It3UsAvTwpAk-9LUnmiz7dRqJnpFX/edit?usp=sharing&ouid=101248488395326982475&rtpof=true&sd=true}{1. Chord learning and adversarial framework for symbolic music generation}}
        \begin{overlayarea}{\textwidth}{0.7\textheight}
            \only<1>{
                \begin{itemize}
                    \item Proposed a chord learning framework for multi-track symbolic music generation based on VAE and GAN.
                    \item Introduced WGAN-GP to solve mode collapse and vanishing gradient problem.
                \end{itemize}
            
                \begin{figure}
                    \centering
                    \includegraphics[scale=0.4]{img/srtp-arch.png}
                \end{figure}
            }
            
            \only<2>{
                \setbeamertemplate{itemize item}{$\checkmark$}
                \begin{itemize}
                    \item CCCC-Mobile Application Innovation Contest\\
                    \textit{First Prize and Most Innovative Award}
                    \item The 4th Internet+ Innovation and Entrepreneurship Competition\\
                    \textit{Gold Award}
                \end{itemize}
            }
        \end{overlayarea}
    \end{block}
\end{frame}

\begin{frame}[t]
    \frametitle{Projects}
    \begin{block}{\href{https://drive.google.com/file/d/1QqgPoRGhaSeS0QWrjnR8mxhvV_onwZ5y/view?usp=drive_link}{2. (Bachelor's Thesis) Voice Conversion Based on Generative Adversarial Networks}}
        \begin{itemize}
            \item Improved StarGAN-VC2 based on multi-speaker non-parallel corpus.
            \item Introduced WGAN-div, AdaIN layer, and neural vocoder, WaveGlow, for better results.
        \end{itemize}

        \begin{figure}[H]
            \centering
            \includegraphics[scale=0.15]{img/undergraduate-thesis-arch.png}
        \end{figure}
    \end{block}
\end{frame}

\begin{frame}[t]
    \frametitle{Projects}
    \begin{block}{3. (Master's Thesis) Floating image quality compensation algorithm technology}
        \begin{itemize}
            \item Proposed a two-stage method to enhanced image quality of floating images: 
            (1) simulating imaging process by solving QP problem derived from optical principles, and 
            (2) refining light field images iteratively based on the simulation.
        \end{itemize}

        \begin{figure}[H]
            \centering
            \includegraphics[scale=0.15]{img/correction-4.png}
        \end{figure}
    \end{block}
\end{frame}

% \begin{frame}[t]
%     \frametitle{Projects}
%     \begin{block}{3. Synthetic Data Generation using Conditional Normalizing Flows}
%         \begin{itemize}
%             \item Ability to replace original datasets and maintain same data distribution, and also highlight outliers for analysis.
%             \item Built an B/S application to visualize results.
%         \end{itemize}

%         \begin{figure}[H]
%             \centering
%             \includegraphics[scale=0.25]{img/cnf.png}
%         \end{figure}
%     \end{block}
% \end{frame}

% TODO: Add architecture.
% \begin{frame}
%     \frametitle{Projects}
%     \begin{block}{4. (Master's Thesis, in progress) Floating image quality compensation algorithm technology}
%         \begin{itemize}
%             \item Simulated imaging process and enhanced the perceived image quality.
%             \item Ability to simulate on different devices with calculations performed only once per simulation.
%         \end{itemize}

%         \begin{figure}[H]
%             \centering
%             \includegraphics[scale=0.25]{img/floating-img-arch.png}
%         \end{figure}
%     \end{block}
% \end{frame}

%----------------------------------------------------------------------------------------
%	AWARDS SLIDES
%----------------------------------------------------------------------------------------

% \section{Awards}
% \begin{frame}
%     \frametitle{Awards}
%     \begin{block}{CCCC-Mobile Application Innovation Contest}
%         First Prize \hfill 09 2018
%     \end{block}
%     \begin{block}{The 4th ”Internet+” Innovation and Entrepreneurship Competition}
%         Gold Award \hfill 09 2018
%     \end{block}
%     \begin{block}{College Students’ Innovative Entrepreneurial Training Plan Program (SRTP)}
%         Excellent \hfill 04 2018 – 05 2019
%     \end{block}
% \end{frame}

%----------------------------------------------------------------------------------------
%	EXPERIENCE SLIDES
%----------------------------------------------------------------------------------------

% \section{Experience}
% \begin{frame}
%     \frametitle{Experience}
%     \begin{block}{The University of British Columbia Vancouver Summer Program}
%         Vancouver, Canada \hfill 07 2019 – 08 2019
%         \begin{itemize}
%             \item Included 2 courses, ”Algorithms and the World Wide Web” and ”Building Modern Web Applications”.
%             \item Collaborated with my teammates to complete our projects as team leader.
%         \end{itemize}
%     \end{block}
% \end{frame}

%----------------------------------------------------------------------------------------
%	SKILLS SLIDES
%----------------------------------------------------------------------------------------

\section{Skills}
\begin{frame}
    \frametitle{Skills}
    \begin{block}{Languages}
        Python, C, C++, Java, JavaScript
    \end{block}
    \begin{block}{Technologies/Frameworks}
        Flask, PyTorch, Spring Boot, Vue.js
    \end{block}
\end{frame}

\section{Thanks}
\begin{frame}
    \begin{center}
        \huge{\texttt{Thanks for your attention}} \\
        \Huge{Tzu-Chun Hsu} \\~\\
        \scriptsize{
            \href{tel:+886-987605719}{ \raisebox{-0.1\height}\faPhone\ \underline{+886-987605719} ~} 
            \href{mailto:vm3y3rmp40719@gmail.com}{\raisebox{-0.2\height}\faEnvelope\  \underline{vm3y3rmp40719@gmail.com}} \\~\\
            \href{https://www.linkedin.com/in/tzu-chun-hsu-ab4b3b188/}{\raisebox{-0.2\height}\faLinkedinSquare\ \underline{tzu-chun-hsu-ab4b3b188} ~}
            \href{https://github.com/Oscarshu0719}{\raisebox{-0.2\height}\faGithub\ \underline{Oscarshu0719}}
        }
    \end{center}
\end{frame}
\end{document}
